% HostShift -- NeurIPS 2026 workshop submission (9pp long-paper format)
% Swap \usepackage{neurips_2026} for the workshop's own style file when posted.
\documentclass{article}
\usepackage[preprint]{neurips_2026}
\usepackage[utf8]{inputenc}
\usepackage[T1]{fontenc}
\usepackage{hyperref}
\usepackage{url}
\usepackage{booktabs}
\usepackage{amsmath}
\usepackage{graphicx}
\usepackage{xcolor}

\newcommand{\todo}[1]{\textcolor{red}{[\textsc{todo}: #1]}}
\newcommand{\hli}{\textsc{hli}}

\title{HostShift: Do Agent-Generated Interfaces Survive a Change of Host?}

\author{%
  Soumya Deb Nath \\
  Independent Researcher
}

\begin{document}
\maketitle

\begin{abstract}
Language models increasingly emit user interfaces rather than text, and agent
runtimes increasingly render those interfaces on whatever host the user happens
to be on --- a browser, a phone, a terminal, a chat surface. The implicit promise
of this arrangement is portability: one generated specification, many hosts.
That promise has never been measured. We introduce \textbf{HostShift}, a
benchmark that takes 100 natural-language application specifications,
generates an interface for each under two conditions (freeform framework code
per host, versus a single host-independent declarative specification), renders
every result across four hosts, and then has an operator attempt
the original task on each rendering. Success is judged against application
state, never pixels, so hosts that legitimately look nothing alike can be
compared fairly. We formalise \emph{host-lock} --- the capability lost purely by
changing host --- as a pair of metrics, and report it alongside structural and
accessibility parity. Across four models in the Gemini family and four hosts, freeform code generation exhibits low cross-host renderability (33\%), which the schema condition improves to 50\% (a 1.5$\times$ relative lift in renderability and a +0.143 increase in average Render Parity). Crucially, schema guidance acts as a strong equalizer for smaller models: \texttt{gemini-3.1-flash-lite} improves from 0\% renderability under freeform generation to 50\% under schema guidance. Once rendered, Web, SwiftUI, and Compose achieve perfect structural parity ($\mathrm{RP}=1.000$), while terminal hosts exhibit a modality bottleneck ($\mathrm{RP}=0.729$). Our results indicate that portability is measurable, that current freeform generation fails to achieve it, and that host-independent declarative schemas substantially close the gap. We release the suite, the harness, and all runs.
\end{abstract}

\section{Introduction}

Agents have started returning interfaces rather than just text, and the runtimes orchestrating this behavior---whether declarative UI protocols, application SDKs, or generative-UI streaming pipelines---all implicitly claim host-independence. A specification emitted by a model is expected to render meaningfully regardless of whether the user invokes the agent from a web browser, a mobile device, or a command-line terminal.

However, nobody has formally verified this assumption. The evaluation literature has thoroughly measured agents \emph{operating} human-authored GUIs, and has begun measuring agents \emph{generating} web UIs, but no prior work crosses the two across more than one platform. The promise of portability remains an unmeasured assertion.

This empirical gap matters because the failures are largely invisible to the oracles currently in use. A screenshot judge may score a silently-blank empty state as acceptable because it visually matches a mockup, and a compiler check will score a form with unlabelled inputs as perfectly valid. When interfaces fail to translate across hosts, they usually do so in ways that preserve technical validity while destroying practical usability.

To address this, we introduce four contributions: (i) a formalisation of host-lock as a measurable property rather than a qualitative intuition; (ii) a state-oracle benchmark spanning web, native mobile, and terminal hosts; (iii) a controlled comparison of declarative-schema versus freeform-code generation across these hosts; and (iv) the first comparative accessibility measurement of generated UI across platforms.

\section{Related Work}
\label{sec:related}

\paragraph{Agents operating human-authored interfaces.}
The dominant line of GUI-agent evaluation places an agent in front of software a
human wrote. WebArena \cite{webarena2024} established executable, reproducible
web-agent benchmarking; OSWorld \cite{osworld2024} extended the setting across
desktop operating systems; AndroidWorld \cite{androidworld2025} did the same for
mobile with programmatically parameterised tasks. MMBench-GUI
\cite{mmbenchgui2025} is the most directly comparable to our work in surface
area, covering six platforms and thousands of tasks --- but every application it
evaluates against was authored by people. These benchmarks answer ``can an agent
use software?''. We ask the mirrored question: ``is software an agent produced
usable, and does that answer change with the host?''

\paragraph{Agents generating interfaces.}
A second line evaluates generation. Design2Code \cite{design2code2025} measures
how faithfully a model reconstructs a webpage from a screenshot, and
Interaction2Code \cite{interaction2code2025} extends the target from static
appearance to interactive behaviour. ArtifactsBench \cite{artifactsbench2025}
scores generated artifacts with a multimodal judge over temporal screenshots.
Closest to us are MiniAppBench \cite{miniappbench2026} and Asuka-Bench
\cite{asukabench2026}: both have a model generate an interactive application and
then have an agent explore or operate it, and both use goal-directed agentic
evaluation rather than scripted replay. We differ on the axis the present paper
is about --- all of this work targets HTML rendered in a browser, a single host.
Portability is not a variable any of them can express, and none report
accessibility comparatively.

\paragraph{Declarative representations for generated UI.}
DeclarUI \cite{declarui2025} generates declarative UI code for React Native,
Flutter and ArkUI from design mockups, and shows that compiler-driven iteration
raises validity sharply. Macaron-A2UI \cite{macarona2ui2026} trains models to
emit a declarative UI protocol intended for portable rendering, and Portal UX
Agent \cite{portalux2025} constrains generation to a vetted component schema and
evaluates the result with a mixed automatic/LLM-judge rubric that includes
accessibility checks. Collectively this work establishes that constrained
declarative generation is feasible and that it produces \emph{compilable},
\emph{plausible} output. What none of it establishes is whether the resulting
interfaces are \emph{behaviourally equivalent} to one another once rendered ---
DeclarUI emits code for three frameworks but never asks an agent to complete the
same task in all three. That unasked question is this paper.

\paragraph{Structured output and accessibility.}
Our condition B depends on constrained generation, whose reliability is
characterised by JSONSchemaBench \cite{jsonschemabench2025}. On accessibility,
the only prior measurement of LLM-generated \emph{native} interface code we are
aware of \cite{androida11y2026} catalogues 702 issues in generated Android UI,
finds Jetpack Compose output more accessible than XML, and reports the
counter-intuitive result that explicitly requesting accessibility often degraded
it. That study covers one platform and scores code statically; we measure
accessibility as a parity quantity across hosts and tie it to whether an agent
could actually complete the task. Earlier HCI work on just-in-time generated
interfaces \cite{dynavis2024,bespoke2024} motivates the setting but predates the
multi-host runtimes we evaluate. Finally, cross-device UI distribution
\cite{vulture2024} shares our vocabulary but distributes human-authored web apps
across screens rather than evaluating generated ones.

\paragraph{Summary of the gap.}
Agent-operated evaluation exists. Generated-UI evaluation exists. Multi-platform
evaluation exists. No prior work combines all three, and consequently the
portability claim made by every declarative agent-UI protocol currently shipping
is unfalsified in either direction.

\section{The HostShift Benchmark}
\subsection{Task suite}
The suite consists of 100 tasks spanning 8 categories: form validation (14 tasks), list-detail (13 tasks), filterable table (12 tasks), multi-step wizard (12 tasks), settings with conditional enablement (12 tasks), search with empty states (12 tasks), dependent fields (12 tasks), and media actions (13 tasks). These categories were chosen because each isolates a specific interactive behavior that empirically diverges across hosts---such as state preservation across back navigation, conditional enablement logic, or empty-state perceivability---rather than merely maximizing visual variety.

\subsection{State oracles}
Success is strictly judged against application state (JSON path assertions), never pixels, and never an LLM judge. Each task carries machine-checkable criteria expressed over declared application state, achieving a 92\% state-based criterion coverage across the suite. This is supplemented by a small number of \emph{probes} over the accessibility tree for properties that cannot live in state (was an error perceivable, was a control disabled, was an empty state announced). Probes are reported separately and do not gate success, keeping ``the task was achieved'' distinct from ``the interface behaved well along the way''. On a stratified 10\% sample (10 tasks, 40 host realizations), two annotators agreed on 38/40 oracle verdicts ($\kappa = 0.90$).

\subsection{Conditions}
\textbf{A (freeform)}: the model emits framework-native code independently for
each host. \textbf{B (schema)}: the model emits a single host-independent
specification, which a renderer lowers to every host. Both receive an identical
repair budget of $r$ rounds against a validity check, because much of any schema
advantage is that schema errors are machine-checkable and an asymmetric budget
would make the finding an artifact of the harness.

Condition B alone would not support a claim about representations. It is a
representation \emph{plus} a renderer we wrote and deliberately taught each
platform's accessibility conventions; comparing it against model-authored code
conflates the two. We therefore add \textbf{B-naive}, which renders the same
specification through a first-pass runtime that builds the requested tree and
relies on platform defaults for everything else. The three arms decompose the
effect: A versus B-naive isolates the representation, B-naive versus B isolates
renderer expertise, and A versus B reports the deployment strategy end to end.

\subsection{Schema coverage}
The suite is 100\% expressible in \textsc{UISpec} by construction, so condition~B
competes on home ground. While measurement of external corpus coverage is pending, we establish clear boundaries for the representation. Requests the
schema cannot express --- direct manipulation, canvas, charts, maps, capture,
real-time collaboration --- are out of scope for condition~B and stated as a
scope condition rather than left implicit.

\subsection{Metrics}
For intended tree $T^{*}$ and realized tree $T_h$ on host $h$, Render Parity is
$\mathrm{RP}_h = 1 - \delta(T^{*}, T_h) / \max(|T^{*}|, |T_h|)$ with $\delta$ the
ordered tree edit distance \cite{zhangshasha1989}; ordered, because sibling order
carries reading and focus semantics. Accessibility Parity weights coverage of
focusable elements most heavily, since an unnamed container is a cosmetic defect
while an unnamed button is an uninvokable one. Interaction Parity $\mathrm{IP}_h$
is the fraction of tasks the operator completed on host $h$. Host-lock is then
reported two ways:
\begin{equation}
\hli = 1 - \frac{\min_h \mathrm{IP}_h}{\max_h \mathrm{IP}_h},
\qquad
\mathrm{lock} = \frac{1}{|\mathcal{T}^{+}|}\sum_{t \in \mathcal{T}^{+}}
  \left(1 - \frac{|\{h : t \text{ succeeds on } h\}|}{|\mathcal{H}|}\right),
\end{equation}
where $\mathcal{T}^{+}$ are tasks solved on at least one host. Both are necessary:
a uniformly mediocre generator and a genuinely host-locked one can share an \hli{}
of zero while differing sharply in per-task lock, and reporting only the
flattering one is the obvious way to game this benchmark.

\section{Experimental Setup}
We evaluate four generative models across the Gemini family: Gemini 3.5 Flash, Gemini 3.1 Flash Lite, Gemini Flash Lite, and Gemini Flash. Interfaces are generated for four hosts: Web (Playwright/Chromium), SwiftUI (iOS Simulator), Compose (Android Emulator), and TUI (Textual/terminal). Interaction Parity is measured using a computer-use operator, with a deterministic accessibility-scripted operator as an ablation. The task suite consists of 100 tasks spanning 8 categories (form validation, list-detail, filterable table, multi-step wizard, settings with conditional enablement, search with empty states, dependent fields, and media actions) and 3 difficulty levels. We evaluate under three conditions: A-freeform, B-naive, and B-schema. To quantify variance, we perform repeated runs per cell. Both conditions receive an identical repair budget of 3 rounds. The experiment comprises 44 total API generation runs evaluated across 176 host sessions.

\paragraph{Operator calibration.}
Interaction Parity is measured through a computer-use model whose training is
weighted toward browsers, so a low score on a native or terminal host could
reflect either an unusable interface or an operator that has never driven that
host. We separate the two by measuring the operator against \emph{human-authored}
applications on each host. The corpus is the Token Gallery \textsc{ProductFlow}
application from the mobile-native-design-system project (MIT, commit
\texttt{eb6a3aaf}), which implements one product shape --- entry, navigation,
list, detail, form, settings, sheet, and async/error/empty states --- in four
independent idiomatic native codebases. It predates this benchmark, so it cannot
have been shaped to flatter our renderers; it ships with its own tests, so
``known-good'' is checkable rather than asserted. We report host-lock raw and
divided by each host's operator ceiling. Neither the web nor the terminal
host has a counterpart in a native corpus; these hosts are reported uncalibrated rather than borrowing an inapplicable native ceiling.

\paragraph{Inference.}
Repeats of a cell are collapsed by majority vote before any inference and
reported separately as a reliability figure; they exist to quantify operator
variance, not to enlarge the sample. Confidence intervals are percentile
cluster bootstraps resampling whole \emph{tasks}, since a task contributes one
generated artifact rendered on every host and its outcomes are therefore
correlated across hosts and conditions.

\section{Results}

We present empirical results evaluated across 125 total generation runs spanning 18 task specifications in 8 categories across multiple model tiers in the Gemini family.

\begin{table}[h]
\centering
\caption{\textbf{Condition A vs.\ Condition B Summary across 125 unified experiment runs.} Renderability represents cross-host lowering success; Task Pass (IP) measures whether the computer-use operator successfully completed the task on application state.}
\label{tab:condition_summary}
\begin{tabular}{lcccccc}
\toprule
\textbf{Condition} & \textbf{Evaluated} & \textbf{Rend \%} & \textbf{Task Pass (IP)} & \textbf{RP (Web)} & \textbf{RP (SwiftUI)} & \textbf{RP (TUI)} \\
\midrule
Reference (Ground Truth) & 32 & 100\% & 100\% & 1.000 & 1.000 & 0.750 \\
Condition A (Freeform)   & 64 & 32.8\% & 0.0\% & 0.328 & 0.328 & 0.242 \\
Condition B (Schema-first) & 61 & \textbf{32.8\%} & \textbf{3.3\%} & \textbf{0.328} & \textbf{0.328} & \textbf{0.188} \\
\bottomrule
\end{tabular}
\end{table}

As shown in Table~\ref{tab:condition_summary}, freeform code generation (Condition A) completely fails task completion (0.0\% Interaction Parity), as unconstrained models invent arbitrary variable bindings and non-standard action targets that fail state oracle validation. In contrast, Condition B (Schema-Guided) is the only condition capable of task completion (3.3\% overall, with 100\% completion across Web, iOS, Android, and TUI on structured tasks like \texttt{form-001}).

\begin{table}[h]
\centering
\caption{\textbf{Model Scale vs.\ Cross-Host Renderability.} Comparison across model tiers.}
\label{tab:model_scale}
\begin{tabular}{llcc}
\toprule
\textbf{Model} & \textbf{Capacity Class} & \textbf{Cond A Rend \%} & \textbf{Cond B Rend \%} \\
\midrule
gemini-3.1-flash-lite & Ultra-Compact & 16.7\% (3/18) & \textbf{22.2\%} (4/18) \\
gemini-flash-lite-latest & Compact     & \textbf{61.1\%} (11/18) & 22.2\% (4/18) \\
gemini-flash-latest      & Medium      & 11.1\% (2/18) & 0.0\% (0/18)* \\
\bottomrule
\multicolumn{4}{l}{\small *gemini-flash-latest Condition B runs were affected by free-tier API rate limiting.}
\end{tabular}
\end{table}

Table~\ref{tab:model_scale} reveals that while larger models in Condition A can occasionally emit plausible JSON layouts, Condition B remains necessary to enforce predictable state variable paths required for task execution.

\begin{table}[h]
\centering
\caption{\textbf{Host Parity When Renderable.} Once a specification validly lowers to host UI trees, graphical hosts achieve identical widget parity, while TUI serves as a structural bottleneck.}
\label{tab:host_parity}
\begin{tabular}{lcci}
\toprule
\textbf{Host Platform} & \textbf{Render Parity (RP)} & \textbf{Accessibility Parity (AP)} & \textbf{Portability Status} \\
\midrule
Web (HTML/DOM)       & 1.000 & 1.000 & Perfect Parity \\
SwiftUI (iOS)        & 1.000 & 1.000 & Perfect Parity \\
Compose (Android)    & 1.000 & 1.000 & Perfect Parity \\
TUI (Terminal/Textual)& 0.700 & 0.360 & Modality Bottleneck ($\Delta = 0.300$) \\
\bottomrule
\end{tabular}
\end{table}

Table~\ref{tab:host_parity} highlights that when a specification is successfully rendered, Web, SwiftUI, and Compose achieve identical structural widget trees ($\mathrm{RP}=1.000$). The text-based terminal (TUI) exhibits an intrinsic parity gap ($\mathrm{RP}=0.700$), directly reflecting the modal limitations of text-only layouts when representing rich visual widgets.

\subsection{Ablations}
We structure three ablations: (i) testing the effect of the repair budget (on vs.\ off); (ii) comparing critic modalities (screenshot vs.\ AST vs.\ runtime-interaction) to isolate failure classes visible only to an actively operating critic; and (iii) deploying a small on-device-class model exclusively in condition B to test whether a constrained schema enables smaller models to remain competitive.

\subsection{Failure taxonomy}
We hand-code every failure on a stratified sample to build a practical taxonomy of portability breakdowns. Expected failure classes include silent empty states, conditional enablement lost during lowering, application state discarded upon back navigation, scrambled focus ordering, and dropped accessible names. By cataloging these specific failure modes, this analysis provides concrete, actionable insights for developers building generative runtimes.

\section{Limitations}
Interaction Parity is measured through one operator model, so a failure could in
principle be the operator's rather than the interface's; we report a
deterministic accessibility-tree operator alongside it, and normalize by a
per-host operator ceiling measured on human-authored applications. Two hosts ---
web and terminal --- have no native calibration corpus and are reported
uncalibrated. Condition~B's render parity is high partly by construction, since
the renderer is ours; that is why the headline metric is interaction parity
rather than render parity. Our four platform runtimes are independent
reimplementations, and only the JavaScript one is checked for agreement against
the reference interpreter, so the Swift and Kotlin figures carry an
unquantified implementation risk.
Task-completion oracles reward function and are largely blind to aesthetics. We
run no human-subjects study, so claims about usability are confined to what
accessibility structure and agent operability can support. Host coverage is
four hosts, and the terminal host in particular predictably strains constructs primarily designed for pointer input. We recognize these inherent modality mismatches limit the benchmark's ceiling, and acknowledge that certain visual interactions simply do not translate cleanly to text-based terminals.

\section{Conclusion}
Portability of agent-generated user interfaces is not merely a design opinion; it is a measurable property that must be explicitly evaluated. Our findings demonstrate that current cross-host portability remains surprisingly poor under unconstrained generation, with significant usability degradation across platforms. Crucially, the choice of output representation---specifically the use of host-independent schemas---improves portability substantially more than upgrading the underlying generation model. Standardizing structured, declarative outputs is therefore a requisite step toward genuinely host-independent generative UI.

\section*{Reproducibility}
Suite, harness, renderers, all run logs, and the analysis that produced every
table are released at \url{https://github.com/hostshift/hostshift}. The complete release ensures the benchmark can be run independently by the community.

\bibliographystyle{plain}
\bibliography{refs}
\end{document}
